\chapter{Welcome \& Vision}
\label{ch:welcome}

\section{What OLK is}

\textbf{Open Library Kashmir (OLK)} was conceived as a community-driven book-sharing initiative to serve readers across the Kashmir Valley. The original operating model relied on local volunteers to collect donated books and redistribute them to students in need. However, the organization rapidly encountered significant operational constraints—including a persistent shortage of volunteers and donors, as well as a limited distribution network that prevented us from reaching our target student demographic. These challenges necessitated a comprehensive strategic redesign of the platform.

The current model introduces a highly efficient, decentralized approach by eliminating nearly all intermediary layers. It establishes a direct, peer-to-peer connection between individuals looking to share their literature and those actively seeking it. Users with surplus books can catalog their collections on the platform, while prospective readers can browse the available inventory and submit a request. The involved parties then coordinate a mutually convenient handover or an agreed-upon shipping arrangement. By removing centralized warehousing and logistical intermediaries, the entire transaction takes place through direct, transparent user-to-user communication.

At the core of this architectural shift is a commitment to fostering meaningful human connections through the shared love of reading. This direct engagement extends beyond the simple exchange of books, enabling a richer, more substantive interaction between users. The Clubs feature is a natural extension of this philosophy—it empowers passionate individuals to create interest-based communities, invite like-minded members, and organize events, all within a single, unified digital ecosystem.

\section{Core features}

\begin{itemize}
  \item \textbf{Community book sharing} — list a book you own as either a \emph{loan} (you get it back) or a \emph{donation} (you don't), with condition, genre, and now (as of the two most recent migrations) a description and publication year enriched automatically from an ISBN scan.
  \item \textbf{Location-aware discovery} — browsing sorts by distance from the viewer using the \code{get\_books\_nearby} RPC (Chapter~\ref{ch:functions}), not just recency.
  \item \textbf{Requesting \& the exchange lifecycle} — a full state machine from \code{pending} through \code{accepted}, \code{handed\_over}, and \code{returned}, tracked per-request with reading-progress and due-date awareness (Chapter~\ref{ch:lifecycle}).
  \item \textbf{Messaging} — a realtime 1:1 chat scoped to each book request, so conversation history is naturally organised by exchange rather than by contact.
  \item \textbf{Wishlists with fuzzy matching} — ask for a title you can't find; the moment anyone lists something close to it, a \code{pg\_trgm}-powered match fires a notification.
  \item \textbf{Reading clubs} — geographically-scoped community groups with posts, membership, and an eligibility gate tied to a user's exchange history.
  \item \textbf{Ratings \& reputation} — a lightweight 1--5 star system that feeds a "Trusted Sharer" badge on public profiles.
  \item \textbf{Reporting \& moderation} — every user and book can be reported; a full admin/moderator subsystem triages, bans, warns, and audits (Chapter~\ref{ch:admin}).
  \item \textbf{Notifications \& email digests} — in-app notifications with a Resend-powered weekly email digest for subscribed users.
\end{itemize}

\section{Who this manual is for}

Primarily, the two or three interns who are about to inherit this codebase. You are not expected to have touched Next.js's App Router before, or Supabase, or even necessarily written Postgres row-level security policies — this manual explains all three from the ground up in Parts III and IV. What it does assume is that you can read TypeScript and SQL.

\begin{olkwarn}
This exact version of Next.js is \textbf{not} the Next.js in your training data or muscle memory. The project's own \pathname{AGENTS.md} says so explicitly: \emph{"This version has breaking changes — APIs, conventions, and file structure may all differ from your training data. Read the relevant guide in \texttt{node\_modules/next/dist/docs/} before writing any code."} A concrete example: Next.js 16 renamed \emph{Middleware} to \emph{Proxy}. This project uses the current \pathname{src/proxy.ts} convention. See Chapter~\ref{ch:proxy} for the full story.
\end{olkwarn}

\section{The bigger picture}

OLK wants two things, in this order. First: get the platform live, publicly, most likely on Vercel, serving real readers in Kashmir. Second, and this is the part worth stating explicitly because it will shape architectural decisions for years: \textbf{eventually run the platform's own infrastructure — including Supabase itself, since it is open source — on a local server physically located in Kashmir}, rather than depending indefinitely on a remote, foreign-hosted database.

That second goal is why Chapter~\ref{ch:local-env} is not a throwaway "how to run this on your laptop" section. Learning to run Supabase's full stack under Docker locally — Postgres, GoTrue auth, PostgREST, Realtime, Storage, Studio, all of it — is a rehearsal for the day that same stack runs on a physical machine in-region instead of on a developer's laptop. Every command in that chapter (\code{supabase start}, \code{supabase stop}, \code{supabase db reset}) is a command that will eventually run against real hardware, not a metaphor for it.

\begin{olknote}
As of this edition, the project already keeps a working local Docker/Supabase environment side-by-side with the hosted one, switched with a single environment variable (see Chapter~\ref{ch:local-env}). This exists specifically so that day-to-day development does not have to touch the production database at all.
\end{olknote}

\section{An infrastructure audit, and what we did with it}
\label{sec:infra-audit}

In July 2026, a colleague of the founder reviewed the project and wrote up a full self-hosting migration plan: VPS + PM2 + Caddy in place of Vercel, a self-hosted Supabase Docker stack (Postgres, GoTrue, PostgREST, Realtime, Storage) in place of the managed cloud project, PgBouncer for connection pooling, a self-hosted mail server (Postfix/Mailu) in place of Resend, and system cron in place of Vercel Cron — all sized for a projected \textbf{100,000} users.

\begin{olknote}
The plan's individual pieces were technically sound — PM2+Caddy, PgBouncer, cron, \code{pg\_dump} backups are all conventional, correct choices. The problems were in the framing, not the components:
\begin{itemize}
  \item \textbf{The 100k figure had no basis.} At the time of the audit this project was pre-launch — no \pathname{/messages} page yet, no real users. Designing connection pooling and disaster recovery around a user count you don't have yet is solving a problem before it exists.
  \item \textbf{"Zero code changes" for the Supabase swap is true narrowly, misleading broadly.} \pathname{src/lib/supabase/server.ts} and \pathname{client.ts} do just read environment variables, so the client code claim holds. But self-hosting Supabase means \emph{you} now own Postgres backups, GoTrue config, the Kong gateway, Realtime, and Storage as running services with their own upgrade cadence and CVEs — a real ongoing burden for a single-developer project, not a config change.
  \item \textbf{Self-hosted email is the highest-risk item and the report understated it.} A fresh VPS IP gets throttled or spam-foldered by Gmail/Outlook regardless of correct SPF/DKIM/DMARC — sending reputation is built over months, not configured. For a notifications/digest feature, landing in spam defeats the point. Resend's free tier is cheap insurance against this failure mode; self-hosting email is a false economy at this stage.
\end{itemize}
\end{olknote}

\textbf{Decision (11-07-2026): none of this ships now.} The project stays on Vercel + managed Supabase + Resend until there are more developers and evidence of real load. The first-year planning target is \textbf{10,000} users, not 100k — reasonable for a regional community platform and not a number that stresses the current managed stack.

\begin{olkwarn}
A VPS is \textbf{not} automatically a step toward the data-sovereignty goal stated earlier in this chapter. Renting a generic VPS (Hetzner, DigitalOcean, AWS, ...) just swaps one third party's control (Vercel/Supabase Inc.) for another's — the data still lives on hardware you don't own, in a jurisdiction you didn't choose. It only becomes a genuine step \emph{toward} sovereignty if it's treated as a rehearsal: running the same self-hosted Docker stack already exercised locally (Chapter~\ref{ch:local-env}) on rented infrastructure first, to prove the operational playbook — backups, upgrades, monitoring, TLS — before ever committing to physical hardware in-region. Stopping at "we rent a VPS instead of Vercel" and calling it done would be a lateral move, not progress on the actual goal.
\end{olkwarn}

\textbf{What's parked for later}, revisited once there's more than one developer and/or real signs of load (not before):
\begin{itemize}
  \item VPS-hosted compute (PM2 + Caddy) — as a rehearsal step per the note above, not a cost-cutting move.
  \item Self-hosted Supabase Docker stack, including PgBouncer once connection limits are actually observed rather than assumed.
  \item Self-hosted mail — only if Resend's pricing or limits actually become a problem; deliverability risk needs a real plan first, not just DNS records.
  \item Automated off-site backups (\code{pg\_dump} to self-hosted MinIO or similar) — worth doing earlier than the rest of this list, independent of the hosting question, once there's real user data worth losing.
\end{itemize}

\textbf{What shipped immediately instead}: indexes on the foreign-key and status columns the audit correctly flagged as missing —\\ \pathname{supabase/migrations/20260711093615\_index\_fk\_and\_status\_columns.sql} adds eleven \code{btree} indexes across \code{book\_requests} (\code{book\_id}, \code{requester\_id}, \code{status}), \code{ratings} (\code{request\_id}, \code{rater\_id}, \code{rated\_user\_id}), and \code{reports} (\code{reporter\_id}, \code{reported\_user\_id}, \code{reported\_book\_id}, \code{status}, \code{assigned\_to} — Chapter~\ref{ch:schema}). A cheap, reversible fix with no architectural bet attached, unlike everything else on this list.

\section{A second pass: readiness for launch, and for 30,000--50,000 users}
\label{sec:scaling-audit}

\textbf{Decision (23-07-2026): the planning target moves from 10,000 to 30,000--50,000 users over the following three months.} A second audit, this time asking two concrete questions — "is this ready to actually go live on Vercel?" and "does it survive that much growth?" — found real bugs rather than architectural concerns, and all of them shipped the same day:

\begin{itemize}
  \item \textbf{A second round of missing indexes}: \code{books.owner\_id}, \code{books.status}, \code{messages.request\_id}, \code{club\_members.user\_id}, \code{bookmarks.user\_id} — the first audit's pass (above) covered \code{book\_requests}/\code{ratings}/\code{reports} but missed these five equally hot columns (Chapter~\ref{ch:schema}).
  \item \textbf{\code{get\_books\_nearby} had no \code{LIMIT}} — every call sorted the \emph{entire} \code{books} table by computed distance before returning it. Now takes an optional \code{p\_limit} (default 200). Fixing this taught a real Postgres lesson about function overloading — see the \textbf{KNOWN ISSUE} box in Chapter~\ref{ch:functions}.
  \item \textbf{\pathname{browse/page.tsx}'s plain query had no bound either} — now \code{.limit(200)}, same stopgap as above (\S\ref{sec:phase1} explains why this is a stopgap, not a fix).
  \item \textbf{The weekly digest cron had an N+1}: one \code{auth.admin.getUserById} call plus one \code{resend.emails.send} call per subscriber, sequential, with no timeout budget. Replaced with paginated \code{listUsers()} (builds an id$\to$email map in batches of 1,000) and \code{resend.batch.send()} (100 emails per call), plus an explicit \code{maxDuration} (Chapter~\ref{ch:server-actions}).
  \item \textbf{\pathname{proxy.ts} and \pathname{(dashboard)/layout.tsx} independently repeated the same auth/profile/settings work} — both called \code{auth.getUser()}, both queried \code{profiles}, both queried \code{platform\_settings} (only the Proxy's copy was cached). The Proxy now forwards its already-computed user id, profile fields, and feature flags to the layout via request headers, so the layout does zero Supabase calls of its own (Chapter~\ref{ch:proxy}).
  \item \textbf{A new \pathname{/api/books} route, cached}: the plain "no search, no filter" Browse query — identical for every visitor — is now served from a 20-second in-memory cache instead of hitting Supabase directly on every page load. Deliberately excluded from the Proxy's matcher, since it needs no session at all (Chapter~\ref{ch:proxy}).
  \item Next.js patched \code{16.2.6} $\to$ \code{16.2.11} (Chapter~\ref{ch:techstack}), and \pathname{error.tsx}/\pathname{loading.tsx}/\pathname{not-found.tsx} were added — none existed before.
\end{itemize}

\begin{olknote}
Every item above was verified against the local Docker stack (Chapter~\ref{ch:local-env}), not just read as a diff: \code{supabase db reset} to prove the migrations replay cleanly, a real headless-browser pass through register/login/browse/logout, and a direct test of the cache's TTL behaviour (insert a row in Postgres, confirm the cached response stays stale, wait, confirm it refreshes). Chapter~\ref{ch:playbook}'s "security \& code-quality audit" entry describes this verify-don't-trust habit in general; this pass is a concrete instance of it.
\end{olknote}

\section{The scaling roadmap, phase by phase}
\label{sec:scaling-roadmap}

None of the fixes above make the platform fast at 30--50k users — they remove the ways it would have \emph{broken}. Getting there smoothly is a separate, staged piece of work. The table below is the whole roadmap at a glance; the sections after it go one level deeper on each row. \textbf{Do not start a phase before its trigger condition is actually true} — this is the same discipline as the "100k had no basis" lesson two sections ago, applied going forward instead of in hindsight.

\begin{center}
\begin{tabularx}{\textwidth}{@{}p{2.1cm} p{3.3cm} Y@{}}
\toprule
\textbf{Phase} & \textbf{Start when...} & \textbf{What changes} \\
\midrule
0 --- Done & --- & Bug fixes above: indexes, bounded queries, batched digest, deduped auth, one small cache. \\
1 --- Harden & Now, in parallel with real user growth toward 30--50k & Real pagination, general rate limiting, error monitoring, a proper fix for the geo query, Supabase/Vercel plan upgrades. \\
2 --- Rehearse & There are 2+ developers \emph{and} Phase 1 is live and stable & Run the already-Dockerised Supabase stack (Chapter~\ref{ch:local-env}) on a rented VPS, in parallel with production, as a dry run. \\
3 --- Sovereignty & The VPS rehearsal has run cleanly for a real stretch (months, not days) & Move the same, already-proven stack onto physical hardware in Kashmir. \\
\bottomrule
\end{tabularx}
\end{center}

\subsection{Phase 1 --- harden the managed stack}
\label{sec:phase1}

Still Vercel + hosted Supabase + Resend — nothing here is an infrastructure change, it is code and dashboard configuration.

\textbf{Code/software changes needed:}
\begin{itemize}
  \item \textbf{Real pagination for \pathname{browse/page.tsx}}, replacing the flat \code{.limit(200)} from \S\ref{sec:scaling-audit}. That limit is a silent-truncation risk, not a real fix: once the catalog holds more than 200 available books, books past the cutoff simply stop appearing, with no error and no indication to the user. Range/offset paging (or a "Load more" button) is the actual fix.
  \item \textbf{General HTTP-level rate limiting} — today's only rate limits are two database triggers scoped to \code{messages} and \code{book\_requests} (Chapter~\ref{ch:functions}). Login, signup, and plain \code{/browse} traffic have no limit at all. A small Upstash Redis (or equivalent) sliding-window check inside \pathname{src/proxy.ts}, ahead of the existing auth/feature-flag logic, is the natural place for this.
  \item \textbf{Error monitoring} — there is currently no Sentry, no Vercel Analytics/Speed Insights, nothing. At today's scale that's a gap; at 30--50k it's a blind spot you cannot afford.
  \item \textbf{A real fix for \code{get\_books\_nearby}'s compute cost} — the \code{p\_limit} added in \S\ref{sec:scaling-audit} bounds the \emph{response size}, not the \emph{work}: Postgres still computes the haversine distance for every row matching the status filter before sorting and truncating. A spatial index (PostGIS, or the lighter-weight \code{earthdistance}/\code{cube} extensions with a GiST index) or a coarser geohash/area-bucket pre-filter are the two realistic paths — both are migrations, both belong in Chapter~\ref{ch:functions} once written.
\end{itemize}

\textbf{Server/Supabase changes needed:}
\begin{itemize}
  \item \textbf{Supabase Pro tier, plus a compute add-on sized to observed load} — the free/shared-compute tier this project has run on so far is not sized for 30--50k users regardless of how clean the queries are.
  \item \textbf{Vercel Pro tier} — Hobby's function-timeout and concurrency ceilings, and its non-commercial terms of service, are not a fit for real production traffic at this scale.
  \item \textbf{Raise Supabase's Auth rate limits} (dashboard $\to$ Authentication $\to$ Rate Limits) — \pathname{proxy.ts}'s \code{auth.getUser()} call (Chapter~\ref{ch:proxy}) is a real network round-trip to Supabase's Auth server on \emph{every} matched request, not a local JWT check. The default limits are tuned for a much smaller project than a 30--50k-user one.
  \item \textbf{Confirm Resend's plan covers the volume} — both the per-notification email and the weekly digest (Chapter~\ref{ch:server-actions}) scale with user count; check this before it becomes a silent failure mode.
\end{itemize}

\begin{olktask}
Pick one item from this phase — real pagination for Browse is the most self-contained — and implement it against the local Docker stack, following the same verify-before-trusting discipline as \S\ref{sec:scaling-audit}: \code{supabase db reset}, then a real browser pass, not just \code{tsc}/\code{lint}/\code{build}.
\end{olktask}

\subsection{Phase 2 --- rehearse self-hosting on a rented VPS}

This is the step the \textbf{WARNING} box earlier in this chapter insists on: proving the operational playbook on rented hardware before ever touching physical hardware in Kashmir. Concretely, here is what is \emph{different} from the managed stack once you take this step, service by service:

\begin{itemize}
  \item \textbf{Connection pooling stops being invisible.} On managed Supabase, PostgREST's internal pool to Postgres is Supabase's problem. Self-hosted, you run PgBouncer (or Supabase's own Supavisor) yourself, and you are the one who sizes \code{pool\_size} against Postgres's \code{max\_connections} and notices when it's wrong.
  \item \textbf{Backups become entirely your responsibility.} No more automatic managed snapshots — a \code{pg\_dump} cron job, with the dump copied \emph{off} the same VPS (the "parked for later" bullet earlier in this chapter about off-site backups stops being optional at this point).
  \item \textbf{TLS and routing are now your config, not a platform default.} Caddy (or similar) in front of PM2-managed Next.js and the self-hosted Supabase Docker services (Chapter~\ref{ch:local-env} is literally the same stack, just no longer on a laptop) — Vercel's automatic edge TLS and global network disappear.
  \item \textbf{Monitoring and alerting are now something you configure and watch}, not a dashboard Vercel/Supabase gave you for free.
  \item \textbf{Upgrade cadence and CVEs for Postgres, GoTrue, PostgREST, Realtime, and Storage become your job to track.} This is the single biggest ongoing burden in this whole roadmap, and the reason the first audit (\S\ref{sec:infra-audit}) was right to push back on doing this before there were real users to justify it.
  \item \textbf{Email stays on Resend.} This is the one thing that should \emph{not} change in this phase — a fresh VPS IP has no sending reputation, and losing digest/notification emails to spam folders defeats the entire feature. Revisit only if Resend's own limits become the actual blocker, per the first audit's conclusion.
  \item \textbf{The \pathname{/api/books} caching pattern (Chapter~\ref{ch:proxy}) still works unchanged} — it's plain Next.js code, indifferent to what's behind \code{NEXT\_PUBLIC\_SUPABASE\_URL}. What's worth adding at this phase is a CDN or edge cache in front of Caddy for public GET responses, since Vercel's edge network was doing that invisibly before.
\end{itemize}

\subsection{Phase 3 --- physical hardware in Kashmir}

Everything from Phase 2 carries over unchanged (pooling, backups, TLS, monitoring, upgrades) — this phase adds the concerns that only exist once the hardware is physically yours, in-region:

\begin{itemize}
  \item \textbf{Power reliability} — a UPS at minimum, a generator plan if the budget allows; grid reliability in the region is a real, known constraint, not a hypothetical one.
  \item \textbf{Internet uplink redundancy} — a second ISP or failover link, since a hyperscaler's backbone redundancy is exactly what you're giving up by leaving Vercel/Supabase's cloud.
  \item \textbf{Off-site, out-of-region backup replication becomes non-negotiable}, not just good practice — cloud infrastructure gives you multi-availability-zone redundancy by default; a single physical server does not, so losing the building without a real off-site copy means losing the platform's only copy of user data.
  \item \textbf{Physical security and access control} for the machine itself.
  \item \textbf{Latency for any users outside the region} — a single in-region server may be measurably slower than the global CDN it replaces for anyone connecting from outside Kashmir. Given the project's stated audience (Chapter~\ref{ch:welcome}), this is very likely an acceptable trade for the sovereignty goal — but it should be a decision made explicitly, not a surprise discovered after the move.
\end{itemize}

\begin{olkwarn}
Re-read the \textbf{WARNING} box earlier in this chapter before starting Phase 3: a VPS rehearsal that goes well is what makes this phase a genuine step toward data sovereignty rather than a repeat of the same "we don't control this hardware" problem one layer down. Do not skip Phase 2 to get here faster.
\end{olkwarn}

\section{A tour of the rest of this manual, in one paragraph}

You will set up the project locally with Docker (Part II), learn how the Next.js frontend is organised page by page and component by component (Part III), learn the Supabase/Postgres backend table by table, function by function, and policy by policy (Part IV), then watch those two halves work together across three real user journeys — exchanging a book, matching a wishlist, moderating a report — traced start to finish (Part V). Part VI is what to do once this is live: deploying, the environment variables that matter, a maintenance playbook organised as "if you need to do X, start here," and a candid accounting of what is still broken or unfinished. Let's begin.
